\section{Sujet n\textsuperscript{o}\,20}
\noindent\begin{minipage}{0.5\linewidth}\footnotesize Office du Baccalauréat du Cameroun\end{minipage}%
\hfill\begin{minipage}{0.45\linewidth}\footnotesize\raggedleft Examen : \textbf{Baccalauréat} · Série : \textbf{C/E}\\
Épreuve : \textbf{Mathématiques} · Durée : \textbf{4\,h} · Coef. : \textbf{7}\end{minipage}
\par\smallskip{\color{onbo}\rule{\linewidth}{1pt}}

\subsection*{Partie A — Évaluation des ressources \hfill (15 points)}

\noindent\textbf{Exercice 1.} \hfill\textit{(3 points)}\\
Dans une plantation de bananiers-plantains à Penja, on estime que chaque rejet planté
donne un régime exportable avec la probabilité $0{,}8$, indépendamment des autres. Un
contrôleur examine au hasard $5$ rejets d'une même parcelle. Soit $X$ la variable
aléatoire égale au nombre de rejets, parmi les $5$, donnant un régime exportable.
\begin{enumerate}[label=\arabic*)]
\item Justifier que $X$ suit une loi binomiale dont on précisera les paramètres. \hfill\textit{(0,75 pt)}
\item Calculer $P(X=4)$ et $P(X\ge 4)$ (valeurs arrondies à $10^{-3}$). \hfill\textit{(1,25 pt)}
\item Déterminer l'espérance $E(X)$ et l'écart-type $\sigma(X)$. \hfill\textit{(1 pt)}
\end{enumerate}

\smallskip
\noindent\textbf{Exercice 2.} \hfill\textit{(3 points)}\\
Chaque régime de la plantation reçoit un numéro d'identification $N$ entier naturel.
\begin{enumerate}[label=\arabic*)]
\item Démontrer que pour tout $n\in\N$, $10^{n}\equiv 1\ [9]$. En déduire que $N$ est
divisible par $9$ si et seulement si la somme de ses chiffres l'est. \hfill\textit{(1,25 pt)}
\item Le numéro $N=\overline{4\,3\,a\,7}$ (en base dix, où $a$ est un chiffre) doit être
divisible par $9$. Déterminer la valeur de $a$. \hfill\textit{(0,75 pt)}
\item Déterminer le reste de la division euclidienne de $7^{2026}$ par $9$. \hfill\textit{(1 pt)}
\end{enumerate}

\smallskip
\noindent\textbf{Exercice 3.} \hfill\textit{(4 points)}\\
Soit $f$ la fonction définie sur $]0\,;+\infty[$ par $f(x)=x-\ln x$, de courbe $(\mathcal{C})$.
Pour tout entier $n\ge 1$, on pose $\displaystyle I_n=\int_1^{\mathrm{e}}(\ln x)^{n}\dd x$.
\begin{enumerate}[label=\arabic*)]
\item Étudier les variations de $f$ et montrer que $f(x)\ge 1$ pour tout $x>0$. \hfill\textit{(1,25 pt)}
\item À l'aide d'une intégration par parties, montrer que $I_1=1$. \hfill\textit{(1 pt)}
\item Établir, par intégration par parties, la relation $I_{n+1}=\mathrm{e}-(n+1)I_n$. \hfill\textit{(1 pt)}
\item En déduire la valeur de $I_2$. \hfill\textit{(0,75 pt)}
\end{enumerate}

\smallskip
\noindent\textbf{Exercice 4.} \hfill\textit{(5 points)}\\
Le plan complexe est muni d'un repère orthonormé direct $(O\,;\,\vec u,\vec v)$. On
considère les points $A$, $B$ et $C$ d'affixes respectives $z_A=2$, $z_B=3+i\sqrt3$ et $z_C=1+i\sqrt3$.
\begin{enumerate}[label=\arabic*)]
\item Calculer $\dfrac{z_B-z_A}{z_C-z_A}$ et en déduire la nature du triangle $ABC$. \hfill\textit{(1,5 pt)}
\item Soit $S$ la similitude directe de centre $A$ qui transforme $C$ en $B$. Déterminer
le rapport et l'angle de $S$. \hfill\textit{(1,5 pt)}
\item Donner l'écriture complexe de $S$ sous la forme $z'=az+b$. \hfill\textit{(1 pt)}
\item Déterminer l'affixe du point $D$, image de $B$ par $S$. \hfill\textit{(1 pt)}
\end{enumerate}

\subsection*{Partie B — Évaluation des compétences \hfill (5 points)}

\noindent\textit{Madame Eyenga exploite la plantation de bananiers-plantains de Penja et
te sollicite, en tant qu'élève de Terminale~C, pour trois décisions de campagne. Elle
dispose de \SI{180}{\metre} de grillage pour clôturer une parcelle rectangulaire
\emph{adossée à une rivière} (le côté longeant la rivière n'est pas clôturé). Sa récolte,
de \SI{15}{\tonne} cette première année, augmente ensuite de \SI{6}{\percent} chaque année
grâce à de nouveaux rejets. Enfin, lors du tri à l'expédition, \SI{10}{\percent} des
régimes présentent un défaut ; un client en choisit $5$ au hasard.}

\begin{enumerate}[label=\textbf{Tâche \arabic*.},leftmargin=2.4em]
\item Déterminer les dimensions de la parcelle rectangulaire d'\emph{aire maximale}
qu'elle peut clôturer, ainsi que cette aire. \hfill\textit{(1,5 pt)}
\item Déterminer l'année à partir de laquelle sa récolte annuelle dépassera \SI{25}{\tonne}. \hfill\textit{(1,5 pt)}
\item Déterminer la probabilité qu'\emph{au moins un} des $5$ régimes choisis présente un défaut. \hfill\textit{(1,5 pt)}
\end{enumerate}
\par\smallskip{\footnotesize\itshape Présentation générale : 0,5 point.}

% ════════════════════════════════════════════════════
\corrige{du sujet 20}

\noindent\textbf{Partie A — Exercice 1.}
1) Les $5$ examens sont indépendants, chacun à deux issues (« exportable » de
probabilité $p=0{,}8$ ou non). Donc $X\sim\mathcal{B}(5\,;\,0{,}8)$.
2) $P(X{=}4)=\binom54(0{,}8)^4(0{,}2)=5\times0{,}4096\times0{,}2=0{,}4096\approx0{,}410$.
$P(X{=}5)=(0{,}8)^5=0{,}32768$, d'où $P(X\ge4)=0{,}4096+0{,}32768=0{,}73728\approx0{,}737$.
3) $E(X)=np=5\times0{,}8=4$ ; $\sigma(X)=\sqrt{np(1-p)}=\sqrt{5\times0{,}8\times0{,}2}=\sqrt{0{,}8}\approx0{,}894$.

\noindent\textbf{Exercice 2.}
1) $10\equiv1\ [9]$, donc par récurrence (ou produit) $10^n\equiv1^n\equiv1\ [9]$. Si
$N=\sum_k c_k 10^k$, alors $N\equiv\sum_k c_k\ [9]$ : $N$ et la somme de ses chiffres ont
le même reste modulo $9$, donc $9\mid N\iff 9\mid\sum_k c_k$.
2) $\overline{4\,3\,a\,7}$ : somme des chiffres $=4+3+a+7=14+a$. On veut $9\mid14+a$ avec
$0\le a\le9$ : $14+a=18\Rightarrow a=4$. Donc $a=4$.
3) Modulo $9$ : $7^1\equiv7$, $7^2=49\equiv4$, $7^3\equiv7\times4=28\equiv1\ [9]$. La
période est $3$. Or $2026=3\times675+1$, donc $7^{2026}\equiv7^1\equiv\mathbf{7}\ [9]$.

\noindent\textbf{Exercice 3.}
1) $f'(x)=1-\frac1x=\frac{x-1}{x}$ : négatif sur $]0,1[$, positif sur $]1,+\infty[$. $f$
décroît puis croît, minimum en $x=1$ avec $f(1)=1-\ln1=1$. Donc $f(x)\ge1$ pour tout $x>0$.
2) IPP avec $u=\ln x$, $v'=1$ ($u'=\frac1x$, $v=x$) :
$I_1=\int_1^{\mathrm e}\ln x\dd x=[x\ln x]_1^{\mathrm e}-\int_1^{\mathrm e}1\dd x=(\mathrm e-0)-(\mathrm e-1)=1$.
3) IPP avec $u=(\ln x)^{n+1}$, $v'=1$ ($u'=(n+1)\frac{(\ln x)^n}{x}$, $v=x$) :
$I_{n+1}=[x(\ln x)^{n+1}]_1^{\mathrm e}-\int_1^{\mathrm e}(n+1)(\ln x)^n\dd x=\mathrm e-(n+1)I_n$.
4) Pour $n=1$ : $I_2=\mathrm e-2I_1=\mathrm e-2$.

\noindent\textbf{Exercice 4.}
1) $z_B-z_A=1+i\sqrt3$, $z_C-z_A=-1+i\sqrt3$. Quotient :
$\frac{1+i\sqrt3}{-1+i\sqrt3}=\frac{(1+i\sqrt3)(-1-i\sqrt3)}{(-1)^2+3}=\frac{-1-i\sqrt3-i\sqrt3+3}{4}=\frac{2-2i\sqrt3}{4}=\frac{1-i\sqrt3}{2}$.
Module $=1$ et argument $=-\frac\pi3$. Comme $|z_B-z_A|=|z_C-z_A|=2$ ($AB=AC$) et l'angle
$(\overrightarrow{AC},\overrightarrow{AB})=-\frac\pi3$, le triangle $ABC$ est isocèle en
$A$ avec un angle de $\frac\pi3$ : il est \textbf{équilatéral}.
2) $S$ de centre $A$ envoie $C$ sur $B$ : son rapport est $k=\frac{|z_B-z_A|}{|z_C-z_A|}=\frac22=1$
et son angle $\theta=\arg\frac{z_B-z_A}{z_C-z_A}=-\frac\pi3$. C'est la rotation de centre
$A$ et d'angle $-\frac\pi3$.
3) $z'-z_A=\mathrm e^{-i\pi/3}(z-z_A)$ avec $\mathrm e^{-i\pi/3}=\frac12-i\frac{\sqrt3}{2}$.
Donc $z'=\left(\frac12-i\frac{\sqrt3}{2}\right)z+b$, où $b=z_A-\mathrm e^{-i\pi/3}z_A=2-2\left(\frac12-i\frac{\sqrt3}{2}\right)=1+i\sqrt3$.
Ainsi $z'=\left(\frac12-i\frac{\sqrt3}{2}\right)z+1+i\sqrt3$.
4) $z_D=z_A+\mathrm e^{-i\pi/3}(z_B-z_A)=2+\left(\frac12-i\frac{\sqrt3}{2}\right)(1+i\sqrt3)$.
Or $\left(\frac12-i\frac{\sqrt3}{2}\right)(1+i\sqrt3)=\frac12+i\frac{\sqrt3}{2}-i\frac{\sqrt3}{2}+\frac{3}{2}=2$.
Donc $z_D=2+2=\mathbf{4}$.

\smallskip
\noindent\textbf{Partie B — Situation.}
\emph{Tâche 1.} Largeur $x$ (deux côtés perpendiculaires à la rivière), longueur
$y=180-2x$. Aire $A(x)=x(180-2x)$ ; $A'(x)=180-4x=0$ en $x=45$. Dimensions :
$x=\SI{45}{\metre}$, $y=\SI{90}{\metre}$ ; aire maximale $45\times90=\SI{4050}{\metre\squared}$.

\emph{Tâche 2.} Récolte de l'année $n$ : $u_n=15\times1{,}06^{\,n-1}$. On résout $u_n>25$,
soit $1{,}06^{\,n-1}>\frac53$, d'où $n-1>\frac{\ln(5/3)}{\ln1{,}06}\approx8{,}77$ : à partir
de la \textbf{10\textsuperscript{e} année} ($u_{10}=15\times1{,}06^9\approx\SI{25,34}{\tonne}$).

\emph{Tâche 3.} $P(\text{au moins un défaut})=1-(0{,}9)^5=1-0{,}59049=\mathbf{0{,}40951}$,
soit environ \SI{41}{\percent}.
